\begin{helper}[A replacement-pair missing color is internal on the deleted pair]
Fix a defective color $\delta$ and its replacement data.  If a color
$\alpha$ is missing under $\varphi$ at either replacement vertex
$u^1_\delta$ or $u^2_\delta$, then there is an individual edge $f$ with
$\varphi(f)=\alpha$ and
\[
 G.\operatorname{ends}(f)=
 \{u^1_{\gamma(\delta)},u^2_{\gamma(\delta)}\}.
\]
\end{helper}
\begin{proof}
Write $X=V(T)$,
$U_\gamma=\{u^1_{\gamma(\delta)},u^2_{\gamma(\delta)}\}$,
$Z=X\setminus U_\gamma$, and
$U_\delta=\{u^1_\delta,u^2_\delta\}$.  Let
\[
 \varphi_\delta=\operatorname{swap}_{\gamma(\delta),\delta}
   (\varphi;v^0_\delta),\qquad X_\delta=V(T_\delta).
\]
The exact replacement identity in \noderef{ScheideReplacementData} is
$X_\delta=Z\cup U_\delta$; it also says that both vertices of
$U_\delta$ lie outside $X$.  The ordered layout in
\noderef{ScheideStructuralConfiguration} puts the base $Y$, in particular
$y_0$, in $Z$.  By \noderef{ScheideReplacementMaximum}, the replacement
triple $(T_\delta,e_0,\varphi_\delta)$ is globally maximum.  Thus
\noderef{ScheideGlobalMaximumOrder} makes $X_\delta$ elementary and closed
under $\varphi_\delta$; the same two conclusions hold for $X$ under
$\varphi$ by applying that node to the original globally maximum triple.

Let $\alpha$ be missing at $u^i_\delta$.  The corresponding one of the
exact edges $f^1_\delta,f^2_\delta$ in
\noderef{ScheideReplacementData} is incident there and has color $\delta$,
so the definition \noderef{MissingColors} gives $\alpha\ne\delta$.
The second edge supplied by \noderef{ScheideSymmetricDefectiveInternal} has
color $\gamma(\delta)$ and exact ends $U_\delta$, so the same argument gives
$\alpha\ne\gamma(\delta)$.  Also
$\gamma(\delta)\ne\delta$: the setup
\noderef{ScheideDefectiveSetup} makes $\gamma(\delta)$ missing at
$v^0_\delta$ while the incident identity $f^0_\delta$ has color $\delta$.

The setup also gives $\alpha_0\ne\gamma(\delta)$ and says that $\alpha_0$
is missing at $y_0$.  Moreover $\alpha_0\ne\delta$.  Indeed, defectivity in
\noderef{DefectiveColors} supplies a $\delta$-colored boundary edge of $X$;
closedness of $X$ then makes $\delta$ present at every vertex of $X$, whereas
$\alpha_0$ is missing at $y_0$.  Thus both $\alpha_0$ and $\alpha$ differ
from the two colors switched in forming $\varphi_\delta$.  The pointwise
formula in \noderef{TwoColorComponentSwap} consequently gives, for every
edge identity $q$,
\[
 \varphi(q)=\alpha_0\Longleftrightarrow
 \varphi_\delta(q)=\alpha_0,
 \qquad
 \varphi(q)=\alpha\Longleftrightarrow
 \varphi_\delta(q)=\alpha.                                  \tag{1}
\]
In particular, the two $(\alpha_0,\alpha)$ incidence graphs are identical.

The agreement assertion \noderef{ScheideReplacementPairFixed} applies to
every identity incident with $u^i_\delta$, so $\alpha$ remains missing at
that vertex under $\varphi_\delta$.  Equation (1) likewise shows that
$\alpha_0$ remains missing at $y_0$.  These are distinct vertices of
$X_\delta$, since the first lies outside $X$ and the second lies in $Z$.
Elementariness of $X_\delta$ therefore gives $\alpha\ne\alpha_0$.

Apply \noderef{ScheideStructuralInternalPair} to $\alpha_0$, using its
missing witness $y_0\in Z$ and
$\alpha_0\ne\gamma(\delta)$.  It returns an exact $\alpha_0$-colored edge
joining the two vertices of $U_\gamma$.  Hence both vertices of $U_\gamma$
belong to the $(\alpha_0,\alpha)$ component $K$ through
$u^1_{\gamma(\delta)}$.  Under $\varphi_\delta$, the component $K$ avoids
every vertex of $X_\delta$ and every edge-node incident with $X_\delta$.
Indeed, its starting vertex is outside $X_\delta$, while $\alpha_0$ and
$\alpha$ have the missing witnesses $y_0$ and $u^i_\delta$ in the closed set
$X_\delta$; this is exactly
\noderef{ClosedSetTwoColorComponentAvoidance}.  Equation (1) identifies
this component, including both its vertex nodes and its individual edge
nodes, with the component defined under $\varphi$.

Switch $\alpha_0$ and $\alpha$ on $K$ and call the resulting coloring
$\psi$.  We verify that $(T,e_0,\psi)$ is globally maximum.  First,
\noderef{TwoColorComponentSwapProper} gives properness because
$\alpha_0\ne\alpha$.  The switch interchanges two colored values and fixes
the uncolored value, so the exact-one-uncolored clause of
\noderef{KTriple} is preserved.  Every displayed edge of $T$ has an
endpoint in $Z$: a base edge has its new base vertex there, and a block edge
has a base endpoint there by the exact attachment clause of
\noderef{ScheideStructuralConfiguration}.  The edge-node avoidance above
therefore fixes every displayed edge identity.

It remains to check every Tashkinov witness, rather than merely assert that
the sequence is unchanged.  For a nonfirst base edge, its witness from
\noderef{TashkinovTree} occurs at an earlier base position, hence in
$Y\subseteq Z$.  For a block edge,
\noderef{ScheideStructuralConfiguration} explicitly supplies a witness in
$Y$ for its block color.  These cases exhaust the exact ordered layout.
The vertex and incident-edge avoidance of $K$ fixes every edge incident with
each such witness, hence fixes its full missing-color set.  Thus all clauses
of \noderef{TashkinovTree} hold for $T$ under $\psi$.  The criticality,
chromatic-index, and degree fields of \noderef{KTriple} are unchanged, and
the displayed vertex set still has the cardinality in
\noderef{GlobalMaximumTriple}.  This reconstructs that globally maximum
triple field by field.

Apply \noderef{ScheideGlobalMaximumOrder} to the reconstructed triple.  It
says that $X$ is closed under $\psi$.  We now exclude every
$\alpha$-edge leaving $U_\gamma$.  Such an edge identity lies in $K$, since
both vertices of $U_\gamma$ lie there.  If its other end is in $Z$, then by
(1) it is an $\alpha$-colored boundary edge of the closed set $X_\delta$
under $\varphi_\delta$, contradicting the missing $\alpha$ at
$u^i_\delta$.  If its other end is outside $X$, the second switch changes
its color to $\alpha_0$; it is then a boundary edge of $X$ under $\psi$,
contradicting the missing $\alpha_0$ at $y_0$, whose incident identities
were fixed by component avoidance.  The endpoint decomposition
$X=Z\cup U_\gamma$ leaves only an edge internal to $U_\gamma$.

If no internal $\alpha$-edge existed, the preceding exclusion would make
$\alpha$ missing at both vertices of $U_\gamma$.  They are distinct members
of $X$ by the vertex injectivity clause of \noderef{TreeSequence}, contrary
to the original elementariness of $X$.  Hence an individual edge $f$ exists
with $\varphi(f)=\alpha$ and exact unordered endpoint pair $U_\gamma$, as
claimed.  Distinct parallel identities have throughout remained distinct;
only their exact endpoint pairs and colors were used.
\end{proof}
